% Unit 2: First Order Linear Differential Equations
% Master Unit Reference Guide (Cheat Sheet, Problems, and Solutions)
% Prepared by Staples Education
% Content informed by the local curriculum modules (2.1 to 2.5) and the
% QUIZ_DATA micro-practice and unit-mastery items for Unit 2.

\documentclass[11pt]{article}

\usepackage[margin=1in]{geometry}
\usepackage{amsmath}
\usepackage{amssymb}
\usepackage{xcolor}
\usepackage[most]{tcolorbox}
\usepackage{enumitem}
\usepackage{hyperref}
\hypersetup{colorlinks=true, linkcolor=headerblue, urlcolor=headerblue}

\tcbuselibrary{skins}

% ---------------------------------------------------------------------------
% Unified style-guide palette (see LATEX_STYLE_GUIDE.md)
% ---------------------------------------------------------------------------
\definecolor{headerblue}{HTML}{1C569E}   % overview / structural framing
\definecolor{accentgreen}{HTML}{2E7D32}  % algorithms / methods
\definecolor{accentorange}{HTML}{E27A1E} % formulas / concept takeaways
\definecolor{workpurple}{HTML}{A020F0}   % worked-example tracking (vibrant)

% ---------------------------------------------------------------------------
% Subsection typography: headerblue numbered sub-module titles
% ---------------------------------------------------------------------------
\usepackage{titlesec}
\titleformat{\subsection}
  {\normalfont\large\bfseries\color{headerblue}}{\thesubsection}{1em}{}

% ---------------------------------------------------------------------------
% Five core custom boxes (unbreakable, title={#1} protected)
% ---------------------------------------------------------------------------

% 1. Blue overview capsule for the topics summary.
\newtcolorbox{topicsbox}{
    enhanced,
    colback=headerblue!6!white, colframe=headerblue,
    coltitle=white, fonttitle=\bfseries,
    title=Unit 2 at a Glance: Core Sub-Topics,
    boxrule=0.6pt, arc=2mm,
    left=3mm, right=3mm, top=2mm, bottom=2mm
}

% 2. Orange formula container for standard forms and master formulas.
\newtcolorbox{formulabox}[1][Master Formula]{
    enhanced,
    colback=accentorange!8!white, colframe=accentorange,
    coltitle=white, fonttitle=\bfseries,
    title={#1},
    boxrule=0.6pt, arc=1mm,
    left=3mm, right=3mm, top=2mm, bottom=2mm
}

% 3. Green algorithm block for explicitly laid-out procedures.
\newtcolorbox{methodbox}[1][Method]{
    enhanced,
    colback=accentgreen!6!white, colframe=accentgreen,
    coltitle=white, fonttitle=\bfseries,
    title={#1},
    boxrule=0.6pt, arc=1mm,
    left=3mm, right=3mm, top=2mm, bottom=2mm
}

% 4. Purple west-bordered worked example for step-by-step solutions.
\newtcolorbox{examplebox}[1][Worked Example]{
    enhanced,
    colback=workpurple!4!white, colframe=workpurple,
    coltitle=white, fonttitle=\bfseries,
    title={#1},
    boxrule=0.4pt, sharp corners,
    borderline west={3pt}{0pt}{workpurple},
    left=5mm, right=3mm, top=2mm, bottom=2mm
}

% 5. Orange thin-bordered final concept takeaway.
\newtcolorbox{conceptbox}[1][Key Takeaway]{
    enhanced,
    colback=accentorange!5!white, colframe=accentorange,
    coltitle=white, fonttitle=\bfseries,
    title={#1},
    boxrule=0.3pt, arc=1mm,
    left=3mm, right=3mm, top=2mm, bottom=2mm
}

\setlist[enumerate]{itemsep=4pt, topsep=4pt}

% Number and index sections plus subsections (e.g., 2.1) two levels deep.
\setcounter{secnumdepth}{2}
\setcounter{tocdepth}{2}

\begin{document}

% ---------------------------------------------------------------------------
% Title block
% ---------------------------------------------------------------------------
\begin{center}
    {\LARGE \bfseries Unit 2: First Order Linear Differential Equations}\\[6pt]
    {\large Prepared by Staples Education}\\[4pt]
    {\itshape Master Unit Reference Guide: Cheat Sheet, Problems, and Solutions}
\end{center}

\vspace{0.4em}

% ---------------------------------------------------------------------------
% Table of contents (master guide)
% ---------------------------------------------------------------------------
\tableofcontents

\vspace{0.4em}

% ---------------------------------------------------------------------------
% Topics overview box
% ---------------------------------------------------------------------------
\begin{topicsbox}
This unit assembles the core solving toolkit of first order theory. Each
sub-topic below builds on the last, moving from raw integration techniques
toward genuine physical models.
\begin{itemize}[leftmargin=1.4em, itemsep=3pt]
    \item \textbf{Separable Equations} --- splitting a rate law into a pure
          $x$ side and a pure $y$ side, then integrating each independently.
    \item \textbf{First-Order Linear Equations and Integrating Factors} ---
          the standard form $y' + P(x)y = Q(x)$ and the factor that collapses
          the left side into a single derivative.
    \item \textbf{Initial Value Problems} --- using a known data point to pin
          the arbitrary constant and select one curve from the solution family.
    \item \textbf{Homogeneous Equations and Substitution} --- the
          substitution $y = vx$ that reduces a degree-matched equation to a
          separable one.
    \item \textbf{Applications: Growth, Decay, and Newton's Law of Cooling}
          --- the exponential models that these methods unlock.
\end{itemize}
\end{topicsbox}

\vspace{0.6em}
\hrule
\vspace{0.6em}

% ===========================================================================
% PART I: CONDENSED CHEAT SHEET
% ===========================================================================
\phantomsection
\addcontentsline{toc}{section}{Part I --- Condensed Cheat Sheet}
\section*{Part I --- Condensed Cheat Sheet}
\setcounter{section}{2}   % subsections auto-number as the unit modules 2.1, 2.2, ...

\subsection{Separable Equations}

A first order equation is \emph{separable} when its right side factors into a
function of $x$ alone times a function of $y$ alone --- that single structural
test is the whole game. Once it factors, you divide and multiply until every
$y$ rides with $dy$ and every $x$ rides with $dx$, then integrate both sides.

\begin{formulabox}[Standard Form and Method: Separable]
A separable equation and its solution scheme:
\[
    \frac{dy}{dx} = g(x)\,h(y)
    \quad\Longrightarrow\quad
    \int \frac{1}{h(y)}\,dy = \int g(x)\,dx + C.
\]
One constant $C$ suffices, since the two constants of integration merge into
a single arbitrary constant.
\end{formulabox}

\noindent\textbf{Pitfall.} Dividing by $h(y)$ silently assumes $h(y) \neq 0$.
Any constant value of $y$ that makes $h(y) = 0$ is an equilibrium solution that
this step can quietly discard --- always check it separately.

\begin{examplebox}[Worked Example 2.1: Solve $\frac{dy}{dx} = 6x^{2}y$]
\textbf{Step 1.} Separate the variables:
\[ \frac{1}{y}\,dy = 6x^{2}\,dx. \]
\textbf{Step 2.} Integrate both sides:
\[ \ln|y| = 2x^{3} + C. \]
\textbf{Step 3.} Exponentiate to isolate $y$:
\[ y = e^{2x^{3} + C} = C\,e^{2x^{3}}. \]
The dropped solution $y = 0$ is recovered as the case $C = 0$.
\end{examplebox}

\subsection{First-Order Linear Equations and Integrating Factors}

A first order equation is \emph{linear} when $y$ and $y'$ each appear only to
the first power and never multiply one another. Before reading off anything,
you must normalize so the derivative carries coefficient one --- only then is
$P(x)$ truly the coefficient of $y$.

\begin{formulabox}[Standard Form and Integrating Factor]
Standard form, integrating factor, and solution:
\[
    \frac{dy}{dx} + P(x)\,y = Q(x),
    \qquad
    \mu(x) = e^{\int P(x)\,dx},
\]
\[
    \frac{d}{dx}\big[\mu(x)\,y\big] = \mu(x)\,Q(x)
    \quad\Longrightarrow\quad
    y = \frac{1}{\mu(x)}\left( \int \mu(x)\,Q(x)\,dx + C \right).
\]
\end{formulabox}

\begin{methodbox}[Algorithm: Solving a First-Order Linear Equation]
\begin{enumerate}[leftmargin=1.6em, itemsep=2pt]
    \item Write the equation in standard form $y' + P(x)\,y = Q(x)$, dividing
          through so the derivative carries coefficient one.
    \item Compute the integrating factor $\mu(x) = e^{\int P(x)\,dx}$.
    \item Multiply through by $\mu(x)$; the left side collapses to the single
          derivative $\frac{d}{dx}\big[\mu(x)\,y\big]$.
    \item Integrate both sides, then divide by $\mu(x)$ to isolate $y$.
\end{enumerate}
\end{methodbox}

\noindent The magic of $\mu(x)$ is that multiplying through turns the entire
left side into the exact product-rule derivative $\frac{d}{dx}[\mu y]$ ---
that is the single fact the whole method rests on.

\begin{examplebox}[Worked Example 2.2: Solve $\frac{dy}{dx} + y = e^{x}$]
\textbf{Step 1.} Read off $P(x) = 1$, so the integrating factor is
\[ \mu = e^{\int 1\,dx} = e^{x}. \]
\textbf{Step 2.} Multiply through; the left side collapses to a derivative:
\[ \frac{d}{dx}\big[e^{x}y\big] = e^{x}\cdot e^{x} = e^{2x}. \]
\textbf{Step 3.} Integrate both sides:
\[ e^{x}y = \frac{1}{2}e^{2x} + C. \]
\textbf{Step 4.} Divide by $e^{x}$:
\[ y = \frac{1}{2}e^{x} + C e^{-x}. \]
\end{examplebox}

\subsection{Initial Value Problems}

The general solution of a first order equation is a whole family of curves
indexed by one arbitrary constant. An \emph{initial condition} $y(x_0) = y_0$
selects exactly one member --- it fixes the constant to the value that forces
the curve through the given point.

\begin{formulabox}[General versus Particular Solution]
A first order equation needs exactly \textbf{one} initial condition. The
general solution carries the free constant; the particular solution has it
determined:
\[
    y = F(x) + C \;\;\xrightarrow{\;y(x_0)=y_0\;}\;\; C = y_0 - F(x_0).
\]
\end{formulabox}

\begin{examplebox}[Worked Example 2.3: Solve $\frac{dy}{dx} = 3x^{2}$, $\;y(1)=5$]
\textbf{Step 1.} Integrate for the general solution:
\[ y = x^{3} + C. \]
\textbf{Step 2.} Apply $y(1) = 5$:
\[ 5 = (1)^{3} + C \;\Longrightarrow\; C = 4. \]
\textbf{Step 3.} State the particular solution:
\[ y = x^{3} + 4. \]
\end{examplebox}

\subsection{Homogeneous First-Order Equations and Substitution}

A function $f(x,y)$ is \emph{homogeneous of degree $n$} when scaling both
inputs pulls out a clean power of the scale factor: $f(tx, ty) = t^{n}f(x,y)$.
An equation $M\,dx + N\,dy = 0$ is homogeneous when $M$ and $N$ share the same
degree --- and that shared degree is precisely what lets the ratio
substitution work.

\begin{formulabox}[The Substitution $y = vx$]
For a homogeneous equation, set
\[
    y = vx, \qquad \frac{dy}{dx} = v + x\frac{dv}{dx},
\]
which always reduces the equation to a \emph{separable} one in $v$ and $x$.
After solving for $v$, back-substitute $v = \dfrac{y}{x}$.
\end{formulabox}

\begin{examplebox}[Worked Example 2.4: Solve $\frac{dy}{dx} = \frac{x + 2y}{x}$]
\textbf{Step 1.} Rewrite the right side and substitute $y = vx$:
\[ \frac{dy}{dx} = 1 + 2\frac{y}{x} = 1 + 2v, \qquad v + x\frac{dv}{dx} = 1 + 2v. \]
\textbf{Step 2.} Cancel $v$ and separate:
\[ x\frac{dv}{dx} = 1 + v \;\Longrightarrow\; \frac{1}{1+v}\,dv = \frac{1}{x}\,dx. \]
\textbf{Step 3.} Integrate and exponentiate:
\[ \ln|1+v| = \ln|x| + C \;\Longrightarrow\; 1 + v = Cx. \]
\textbf{Step 4.} Back-substitute $v = \frac{y}{x}$ and clear:
\[ 1 + \frac{y}{x} = Cx \;\Longrightarrow\; y = Cx^{2} - x. \]
\end{examplebox}

\subsection{Applications: Growth, Decay, and Newton's Law of Cooling}

When a rate of change is proportional to the current amount, the solution is
exponential --- this single idea drives population models, radioactive decay,
and cooling alike.

\begin{formulabox}[Exponential Models]
Growth and decay:
\[
    \frac{dP}{dt} = kP \;\Longrightarrow\; P(t) = P_{0}e^{kt}
    \qquad (k>0 \text{ grows},\; k<0 \text{ decays}),
\]
with doubling time $t = \dfrac{\ln 2}{k}$. Newton's law of cooling:
\[
    \frac{dT}{dt} = -k\,(T - T_{s}) \;\Longrightarrow\;
    T(t) = T_{s} + (T_{0} - T_{s})\,e^{-kt},
\]
so $T \to T_{s}$ as $t \to \infty$.
\end{formulabox}

\begin{examplebox}[Worked Example 2.5: A culture doubles in 3 hours]
\textbf{Step 1.} The model $P = P_{0}e^{kt}$ doubles when $e^{3k} = 2$:
\[ 3k = \ln 2 \;\Longrightarrow\; k = \frac{\ln 2}{3}. \]
\textbf{Step 2.} The growth law is therefore
\[ P(t) = P_{0}\,e^{(\ln 2 / 3)\,t} = P_{0}\,2^{\,t/3}. \]
Every $3$ hours the population multiplies by two, as required.
\end{examplebox}

\vspace{0.4em}
\hrule
\vspace{0.6em}

% ===========================================================================
% PART II: REVIEW GUIDE (PRACTICE QUESTIONS)
% ===========================================================================
\clearpage
\phantomsection
\addcontentsline{toc}{section}{Part II --- Review Guide: Practice Problems}
\section*{Part II --- Review Guide: Practice Problems}

Work the following ten problems in order; they are graded from direct
separation through linear factors, initial value problems, the homogeneous
substitution, and a full physical model. Solutions follow in Part III.

\begin{enumerate}[leftmargin=2em]
    \item Solve $\dfrac{dy}{dx} = \dfrac{x}{y}$ by separation of variables.

    \item Solve $\dfrac{dy}{dx} = e^{\,x - y}$.

    \item Solve $\dfrac{dy}{dx} = y\cos x$, and identify any solution that is
          lost when you divide by $y$.

    \item Solve the linear equation $\dfrac{dy}{dx} + 3y = 6$ using an
          integrating factor.

    \item Solve $x\dfrac{dy}{dx} + y = x^{2}$ by first writing it in standard
          linear form.

    \item Solve $\dfrac{dy}{dx} + \dfrac{2}{x}\,y = x$.

    \item Solve the initial value problem $\dfrac{dy}{dx} + 2y = 4$ with
          $y(0) = 1$.

    \item Solve the homogeneous equation $\dfrac{dy}{dx} = \dfrac{x + y}{x}$
          using the substitution $y = vx$.

    \item Solve the homogeneous equation
          $\dfrac{dy}{dx} = \dfrac{x^{2} + y^{2}}{xy}$ using $y = vx$.

    \item \textbf{(Application)} An object at $100^{\circ}$ is placed in a room
          held at $20^{\circ}$. After $10$ minutes its temperature is
          $60^{\circ}$. Find the temperature function $T(t)$ and determine how
          long it takes the object to reach $30^{\circ}$.
\end{enumerate}

\vspace{0.4em}
\hrule
\vspace{0.6em}

% ===========================================================================
% PART III: SOLUTIONS
% ===========================================================================
\clearpage
\phantomsection
\addcontentsline{toc}{section}{Part III --- Complete Solutions}
\section*{Part III --- Complete Solutions}

\noindent\textbf{Solution 1.}\quad Separate, then integrate.
\[ \textcolor{workpurple}{y\,dy = x\,dx} \]
\[ \textcolor{workpurple}{\frac{y^{2}}{2} = \frac{x^{2}}{2} + C_{1}} \]
\[ \textcolor{workpurple}{y^{2} = x^{2} + C} \qquad\text{(general solution).} \]

\medskip
\noindent\textbf{Solution 2.}\quad Split the exponential as
$e^{x-y} = e^{x}e^{-y}$, then separate.
\[ \textcolor{workpurple}{e^{\,y}\,dy = e^{\,x}\,dx} \]
\[ \textcolor{workpurple}{e^{\,y} = e^{\,x} + C} \]
\[ \textcolor{workpurple}{y = \ln\!\left(e^{\,x} + C\right)}. \]

\medskip
\noindent\textbf{Solution 3.}\quad Separate (assuming $y \neq 0$), integrate,
and exponentiate.
\[ \textcolor{workpurple}{\frac{1}{y}\,dy = \cos x\,dx} \]
\[ \textcolor{workpurple}{\ln|y| = \sin x + C} \]
\[ \textcolor{workpurple}{y = C\,e^{\sin x}}. \]
Dividing by $y$ assumed $y \neq 0$, so it dropped the equilibrium solution
$\textcolor{workpurple}{y = 0}$ --- which is recovered as the special case $C = 0$.

\medskip
\noindent\textbf{Solution 4.}\quad Here $P = 3$, so the integrating factor is
$\textcolor{workpurple}{\mu = e^{\int 3\,dx} = e^{3x}}$. Multiply through:
\[ \textcolor{workpurple}{\frac{d}{dx}\big[e^{3x}y\big] = 6\,e^{3x}} \]
\[ \textcolor{workpurple}{e^{3x}y = \int 6\,e^{3x}\,dx = 2\,e^{3x} + C} \]
\[ \textcolor{workpurple}{y = 2 + C\,e^{-3x}}. \]

\medskip
\noindent\textbf{Solution 5.}\quad Divide by $x$ to reach standard form:
\[ \textcolor{workpurple}{\frac{dy}{dx} + \frac{1}{x}\,y = x}, \qquad
   \textcolor{workpurple}{\mu = e^{\int (1/x)\,dx} = e^{\ln|x|} = x}. \]
Multiply through and integrate:
\[ \textcolor{workpurple}{\frac{d}{dx}\big[xy\big] = x\cdot x = x^{2}} \]
\[ \textcolor{workpurple}{xy = \frac{x^{3}}{3} + C} \]
\[ \textcolor{workpurple}{y = \frac{x^{2}}{3} + \frac{C}{x}}. \]

\medskip
\noindent\textbf{Solution 6.}\quad With $P = \frac{2}{x}$, integrate the
coefficient and exponentiate:
\[ \textcolor{workpurple}{\mu = e^{\int (2/x)\,dx} = e^{2\ln|x|} = x^{2}}. \]
Multiply through and integrate:
\[ \textcolor{workpurple}{\frac{d}{dx}\big[x^{2}y\big] = x^{2}\cdot x = x^{3}} \]
\[ \textcolor{workpurple}{x^{2}y = \frac{x^{4}}{4} + C} \]
\[ \textcolor{workpurple}{y = \frac{x^{2}}{4} + \frac{C}{x^{2}}}. \]

\medskip
\noindent\textbf{Solution 7.}\quad The integrating factor is
$\textcolor{workpurple}{\mu = e^{\int 2\,dx} = e^{2x}}$. Multiply through and
integrate:
\[ \textcolor{workpurple}{\frac{d}{dx}\big[e^{2x}y\big] = 4\,e^{2x}} \]
\[ \textcolor{workpurple}{e^{2x}y = 2\,e^{2x} + C} \]
\[ \textcolor{workpurple}{y = 2 + C\,e^{-2x}}. \]
Apply the initial condition $y(0) = 1$:
\[ \textcolor{workpurple}{1 = 2 + C \;\Longrightarrow\; C = -1} \]
\[ \textcolor{workpurple}{y = 2 - e^{-2x}}. \]

\medskip
\noindent\textbf{Solution 8.}\quad Rewrite as $\frac{dy}{dx} = 1 + \frac{y}{x}$
and substitute $y = vx$, so $\frac{dy}{dx} = v + x\frac{dv}{dx}$:
\[ \textcolor{workpurple}{v + x\frac{dv}{dx} = 1 + v} \]
\[ \textcolor{workpurple}{x\frac{dv}{dx} = 1 \;\Longrightarrow\; dv = \frac{1}{x}\,dx} \]
\[ \textcolor{workpurple}{v = \ln|x| + C}. \]
Back-substitute $v = \frac{y}{x}$:
\[ \textcolor{workpurple}{y = x\ln|x| + Cx}. \]

\medskip
\noindent\textbf{Solution 9.}\quad Substitute $y = vx$. The right side becomes
\[ \textcolor{workpurple}{\frac{x^{2} + v^{2}x^{2}}{x\cdot vx} = \frac{1 + v^{2}}{v}}. \]
With $\frac{dy}{dx} = v + x\frac{dv}{dx}$,
\[ \textcolor{workpurple}{v + x\frac{dv}{dx} = \frac{1 + v^{2}}{v} = \frac{1}{v} + v} \]
\[ \textcolor{workpurple}{x\frac{dv}{dx} = \frac{1}{v} \;\Longrightarrow\; v\,dv = \frac{1}{x}\,dx} \]
\[ \textcolor{workpurple}{\frac{v^{2}}{2} = \ln|x| + C_{1} \;\Longrightarrow\; v^{2} = 2\ln|x| + C}. \]
Back-substitute $v = \frac{y}{x}$:
\[ \textcolor{workpurple}{\frac{y^{2}}{x^{2}} = 2\ln|x| + C \;\Longrightarrow\; y^{2} = x^{2}\big(2\ln|x| + C\big)}. \]

\medskip
\noindent\textbf{Solution 10.}\quad Newton's law of cooling with surrounding
temperature $T_{s} = 20$ and initial temperature $T_{0} = 100$ gives
\[ \textcolor{workpurple}{T(t) = 20 + (100 - 20)\,e^{-kt} = 20 + 80\,e^{-kt}}. \]
Use the reading $T(10) = 60$ to find $k$:
\[ \textcolor{workpurple}{60 = 20 + 80\,e^{-10k} \;\Longrightarrow\; 80\,e^{-10k} = 40} \]
\[ \textcolor{workpurple}{e^{-10k} = \tfrac{1}{2} \;\Longrightarrow\; -10k = -\ln 2 \;\Longrightarrow\; k = \frac{\ln 2}{10}}. \]
Hence the temperature function is
\[ \textcolor{workpurple}{T(t) = 20 + 80\,e^{-(\ln 2 / 10)\,t} = 20 + 80\left(\tfrac{1}{2}\right)^{t/10}}. \]
To reach $30^{\circ}$, solve $T(t) = 30$:
\[ \textcolor{workpurple}{30 = 20 + 80\,e^{-kt} \;\Longrightarrow\; 80\,e^{-kt} = 10 \;\Longrightarrow\; e^{-kt} = \tfrac{1}{8}} \]
\[ \textcolor{workpurple}{-kt = -\ln 8 = -3\ln 2 \;\Longrightarrow\; t = \frac{3\ln 2}{k} = \frac{3\ln 2}{\ln 2 / 10} = 30 \text{ minutes}}. \]

\vspace{0.8em}
\begin{conceptbox}[Unit 2 Key Takeaway]
Every method in this unit reduces to a single move --- rewrite the equation
into a form you can integrate. \textcolor{accentorange}{Separation of variables}
splits a factored rate law; the \textcolor{accentorange}{integrating factor
$\mu = e^{\int P\,dx}$} collapses a linear left side into one derivative; and
the \textcolor{accentorange}{substitution $y = vx$} turns a homogeneous equation
separable. Read the structure first, and the technique always follows.
\end{conceptbox}

\end{document}
